\section{Introduction}

An AI agent sends its whole history to a language model, waits for an answer, runs a tool, appends the result and
calls the model again. Each call repeats a long prefix the server has already processed. Serving engines keep the
attention keys and values of that prefix (the KV cache) so that the next call does not recompute them, and the saving
is large: for a 32-billion-parameter model we measure 56--73\,mJ of above-idle GPU energy per recomputed prompt token and
almost nothing per cached one.

The cache has a weakness that agents expose. The engine does not know that an agent has gone to run a tool or how long
it will be away, so it keeps KV by recency and evicts the least recently used entry when memory runs out. Once the agents'
working set exceeds the cache by a factor between 1.2 and 2.2, depending on how long they pause, hits collapse until nearly every agent's KV is evicted before it returns. Prior work has named this collapse
(``middle-phase thrashing'' in CONCUR~\cite{chen2026concur}), located it~\cite{he2026robot,lu2026metastability} and
shown that a host-memory tier helps only when it holds the agents' working set~\cite{shao2026efficientagent}. All of
that work reports latency or throughput. Energy has been measured for KV offloading (at low concurrency, across
buffer sizes of 8--128\,GiB, on production agent traces)~\cite{mao2026offloading} and for prefix reuse on single
requests~\cite{shewale2026prefixbench}, but not through the collapse that many concurrent agents cause, and not for
the admission rules proposed to prevent it.

We report measurements, a small model that predicts them, and an admission gate built from the model. We ran
everything on one NVIDIA H200 NVL with two unmodified engines (vLLM and SGLang) and three models of different shapes (a
dense 32B model, a 30B mixture-of-experts (MoE) model and a dense 8B model); only the 10-minute adaptive-gate runs used a second H200 NVL host. The workloads go from controlled to real:
synthetic sessions whose size and pauses we set, 400 recorded SWE-agent trajectories, a replay of 1/8 and 1/4 samples of two 15-minute windows of
Kimi's one-hour production tool-agent trace, with their original timing and prefix sharing, and 400 graded tasks from the
Berkeley Function Calling Leaderboard (BFCL) multi-turn benchmark. We wrote each prediction down before the run that
tested it (Appendix~\ref{app:prereg}) and kept the ones that failed.

Figure~\ref{fig:arch} shows the paths a prompt token can take and where the gate sits.

\begin{figure*}[t]
\centering
\frenchspacing
\begin{tikzpicture}[
  font=\footnotesize,
  box/.style={draw, rounded corners=2pt, align=center, inner sep=3pt, minimum height=0.68cm,
    execute at begin node={\hyphenpenalty=10000\exhyphenpenalty=10000}},
  dark/.style={box, fill=black, text=white},
  grey/.style={box, fill=black!12},
  light/.style={box, fill=white},
  arr/.style={-{Stealth[length=2mm]}, thick},
  lab/.style={font=\scriptsize\itshape, text=black!60, align=center,
    execute at begin node={\hyphenpenalty=10000\exhyphenpenalty=10000}},
  head/.style={font=\scriptsize\bfseries, align=left}
]
% ---------- agents
\node[head] at (1.3,2.65) {AGENTS ($N$ sessions)};
\node[light, text width=2.4cm] (a1) at (1.3,1.8) {model call $\rightarrow$ tool};
\node[light, text width=2.4cm] (a2) at (1.3,0.9) {model call $\rightarrow$ tool};
\node[light, text width=2.4cm] (a3) at (1.3,0.0) {model call $\rightarrow$ tool};
\node at (1.3,-0.6) {$\vdots$};
\node[light, text width=2.4cm] (an) at (1.3,-1.3) {model call $\rightarrow$ tool};
\node[lab, anchor=north, text width=3.0cm] at ([yshift=-0.12cm]an.south) {each call resends the whole history ($s$~tokens); tool pauses last seconds to minutes};

% ---------- gate
\node[head] at (5.0,2.65) {ADMISSION GATE};
\node[dark, text width=2.3cm] (all) at (5.0,1.8) {admit all $N$};
\node[diamond, draw, aspect=2.2, inner sep=0.5pt, align=center, fill=white, font=\scriptsize] (dec) at (5.0,0.45)
  {$S+N(\bar s-S)\le\theta C$?};
\node[grey, text width=2.3cm] (k) at (5.0,-1.3) {admit $K$ (Eq.~\eqref{eq:gate});\\others wait};
\draw[arr] (dec.north) -- node[left, font=\scriptsize]{yes} (all.south);
\draw[arr] (dec.south) -- node[left, font=\scriptsize]{no} (k.north);
\draw[arr] (a1.east) -- ++(0.35,0) |- (dec.west);
\draw[arr] (an.east) -- ++(0.35,0) |- (dec.west);
\draw[thick] (a2.east) -- ++(0.35,0);
\draw[thick] (a3.east) -- ++(0.35,0);

% ---------- engine
\node[head] at (11.6,2.65) {SERVING ENGINE (vLLM or SGLang) ON ONE H200};
\node[light, text width=1.2cm] (lk) at (8.15,0.45) {prefix\\lookup};
\draw[arr] (all.east) -- ++(0.55,0) |- (lk.west);
\draw[arr] (k.east) -- ++(0.55,0) |- (lk.west);

\node[dark, text width=3.6cm] (hit) at (11.6,1.8) {GPU hit (HBM, capacity $C$)\\[-1pt]{\scriptsize 0--0.3\,mJ per token}};
\node[grey, text width=3.6cm] (res) at (11.6,0.45) {restore from host tier (host memory)\\[-1pt]{\scriptsize 0.27--0.68\,mJ per token}};
\node[light, text width=3.6cm, very thick] (rec) at (11.6,-0.9) {recompute (prefill)\\[-1pt]{\scriptsize 56--73\,mJ per token}};
\draw[arr] (lk.east) -- ++(0.45,0) |- (hit.west);
\draw[arr] (lk.east) -- (res.west);
\draw[arr] (lk.east) -- ++(0.45,0) |- (rec.west);

\node[light, text width=1.3cm] (out) at (15.9,0.45) {decode\\answer};
\draw[arr] (hit.east) -- ++(0.45,0) |- (out.west);
\draw[arr] (res.east) -- (out.west);
\draw[arr] (rec.east) -- ++(0.45,0) |- (out.west);

\node[lab, anchor=north, text width=5.4cm] (meter) at ([yshift=-0.12cm]rec.south) {all three paths metered as GPU board energy (NVML counter, read every 0.1\,s)};

% ---------- cliff callout
\node[draw, dashed, rounded corners=2pt, align=left, font=\scriptsize, text width=8.3cm, fill=black!4] at (11.6,-3.2)
{\textbf{The cliff.} Without the gate, once $W=Ns$ exceeds $\kappa C$ ($\kappa\approx1.2$--2.2), LRU evicts nearly every
session before it returns, so every turn takes the recompute path and costs what it would with no cache.};
\end{tikzpicture}
\caption{What happens to one agent turn. Agents resend their history after each tool call. The admission gate
(Section~\ref{sec:control}) admits every session while the working set fits under $\theta C$ and only $K$ sessions
otherwise ($S$, $\bar s$, $\theta$, $\kappa$ and $K$ are defined in Section~\ref{sec:model}). In the engine, each prompt token takes one of three paths whose above-idle GPU energy we measured
(Section~\ref{sec:cost}): its KV is still in GPU memory, in a host tier, or gone. HBM: the GPU's high-bandwidth memory; LRU: least-recently-used eviction; NVML: NVIDIA Management Library.}
\label{fig:arch}
\end{figure*}


\paragraph{Findings.}
\begin{enumerate}
\item \textbf{Past the cliff, a full cache saves almost no energy.} GPU energy per turn equals that of running with no cache,
to within 0.01--1.6\% in every cell past the cliff we measured (vLLM: three models and two context lengths; SGLang: Qwen3-32B). The cliff falls where the agents'
working set $W$ exceeds the cache capacity $C$ by a factor of 1.2--2.2: hits fall to half their 16-session level at 1.2 and vanish by 1.55 with 5\,s pauses, and fall to
half at 1.5 and vanish by 2.2 with 30\,s pauses. With sessions 2.5 times longer (27k instead of 11k tokens) the cliff moved from between 16 and 32 sessions to between 8 and 16, consistent with the $1/s$ scaling of Eq.~\eqref{eq:cliff} (Section~\ref{sec:cliff}).
\item \textbf{Holding KV costs no measurable power; rebuilding it does.} A full cache changed idle power by less than
0.2\%, and restoring a prompt's KV from host memory took 1/110 to 1/220 of the above-idle GPU energy of recomputing it,
with the gap growing with context length (Section~\ref{sec:cost}).
\item \textbf{On production traffic, a host tier adds capacity.} In a replay of Kimi's tool-agent trace, a 160\,GB host
tier let one H200 finish every request where the GPU cache alone finished 85--86\%, and process 16--18\% more prompt tokens
for about the same GPU energy (15--17\% less per prompt token), in two time windows of a 1/4 sample that overloads the
card; on a 1/8 sample, which the GPU cache served in full, the tier saved under 1\% (Sections~\ref{sec:tiers}
and~\ref{sec:production}).
\item \textbf{The energy is predictable on a busy GPU.} A one-parameter model fitted on synthetic sessions predicts the
production energy ratios to within about 1\% when the GPU is saturated, and a block-level hit-rate predictor run on
hashed request logs matches the measured hit rate to within 0.04--4.3 percentage points when the GPU keeps up. The model fails at
light load and the predictor under overload; we report both (Sections~\ref{sec:model} and~\ref{sec:production}).
\item \textbf{A gate that acts only under pressure costs nothing below the cliff and saves energy above it.} Admitting every agent while the working set is at most 1.2 times the cache, and
$K$ agents from Eq.~\eqref{eq:gate} above that, cost no energy below the cliff, with 90th-percentile resume latency no worse. Above it, in 150\,s
steps, it cut energy per turn by 41--45\% but left a third to a half of the sessions unserved at the highest loads; in
600\,s steps it served every session, completed 48--62\% more turns and cut energy per turn by 42--50\%, at the price of
queueing: the 90th percentile of first-turn latency was 1.9--2.7 times the GPU cache's (Section~\ref{sec:control}).
\item \textbf{On graded tasks, the cache doubled correct answers per kWh.} On BFCL, accuracy with and without the cache differed by 3 of 400 tasks (204 correct with the cache,
201 without; one round) and halved time and energy: 2.0$\times$ the correct answers per kWh (Section~\ref{sec:accuracy}).
\end{enumerate}

\paragraph{What is not new.} Thrashing, the cliff, the weakness of eviction policies, host-tier sizing by working set
and working-set admission control are all published~\cite{chen2026concur,he2026robot,lu2026metastability,
shao2026efficientagent,denning1968working}, and so are offloading during tool calls~\cite{xia2026mori,tokencake2025}
and hit-rate prediction from traces~\cite{kvset2026}. Our contribution is to measure these in GPU energy, on one
hardware platform and one set of workloads, including production traffic; to show how far a simple model predicts them; and to report where
the published ideas fail in practice (Section~\ref{sec:control}). All energy figures are GPU board energy read from
the energy counter of NVML (the NVIDIA Management Library). We could not meter host CPU and memory energy in our rented container, and we do not convert
our percentages to whole-server savings (Section~\ref{sec:discussion}, Limitations).

\begin{figure*}[t]
\centering
\frenchspacing
\begin{tikzpicture}[
  font=\footnotesize,
  box/.style={draw, rounded corners=2pt, align=center, inner sep=3pt, minimum height=0.62cm,
    execute at begin node={\hyphenpenalty=10000\exhyphenpenalty=10000}},
  dark/.style={box, fill=black, text=white},
  grey/.style={box, fill=black!12},
  light/.style={box, fill=white},
  arr/.style={-{Stealth[length=2mm]}, thick},
  head/.style={font=\scriptsize\bfseries, align=left}
]
% ---------- workloads
\node[head] at (1.75,3.3) {WORKLOADS};
\node[light, text width=3.6cm] (w1) at (1.75,2.65) {synthetic sessions ($\sim$11k and $\sim$27k tokens)};
\node[light, text width=3.6cm] (w2) at (1.75,1.8) {400 SWE-agent trajectories};
\node[light, text width=3.6cm] (w3) at (1.75,0.95) {Kimi tool-agent trace (23,608 requests)};
\node[light, text width=3.6cm] (w4) at (1.75,0.1) {BFCL v4 (400 graded tasks)};
\node[grey, text width=2.2cm] (same) at (5.75,1.42) {identical prompts (BFCL: tasks) for every policy};
\foreach \i in {1,2,3,4} \draw[arr] (w\i.east) -- ++(0.35,0) |- (same.west);

% ---------- policies
\node[head] at (9.8,3.3) {POLICIES};
\node[light, text width=4.5cm] (p1) at (9.8,2.65) {no cache; GPU cache};
\node[light, text width=4.5cm] (p2) at (9.8,1.8) {host tier: vLLM 70\,GiB; SGLang 80--300\,GB (WT, WB)};
\node[light, text width=4.5cm] (p3) at (9.8,0.95) {LRU, LFU, SLRU eviction};
\node[light, text width=4.5cm] (p4) at (9.8,0.1) {gates: static, formulas, adaptive};
\foreach \i in {1,2,3,4} \draw[arr] (same.east) -- ++(0.12,0) |- (p\i.west);

% ---------- engine, predictions, measurement
\node[box, dashed, text width=3.6cm] (pre) at (14.75,2.55) {28 predictions, written before the runs};
\node[dark, text width=3.6cm] (eng) at (14.75,1.42) {vLLM 0.33, SGLang 0.5.23\\on an H200 NVL};
\node[grey, text width=3.6cm] (meas) at (14.75,0.3) {GPU joules (NVML), hits, turns, TTFT, accuracy};
\foreach \i in {1,2,3,4} \draw[arr] (p\i.east) -- ++(0.12,0) |- (eng.west);
\draw[arr, dashed] (pre) -- (eng);
\draw[arr] (eng) -- (meas);

% ---------- results
\node[head, anchor=west] at (-0.15,-0.62) {RESULTS};
\draw[thick] (meas.south) -- (14.75,-0.85);
\draw[thick] (1.08,-0.85) -- (15.43,-0.85);
\foreach \x/\s/\t [count=\j] in {1.08/{\ref{sec:cost}}/{cost of a token},
                      3.95/{\ref{sec:cliff}}/{cliff and memory wall},
                      6.82/{\ref{sec:tiers}}/{host tiers},
                      9.69/{\ref{sec:production}}/{production replay},
                      12.56/{\ref{sec:control}}/{admission gate},
                      15.43/{\ref{sec:accuracy}}/{correct per kWh}} {
  \node[box, dashed, text width=2.25cm, minimum height=0.85cm, anchor=north, font=\scriptsize] (r\j) at (\x,-1.15) {\textbf{Section~\s}\strut\\\t\strut};
  \draw[arr] (\x,-0.85) -- (r\j.north);
}
\end{tikzpicture}
\caption{Method. Four workloads, from controlled to real, are sent with identical prompts (for BFCL, identical tasks) to every policy on two
engines and one GPU type\@. Every run records GPU energy from the NVML counter, cache hits, completed turns or requests,
latency and, for BFCL, graded accuracy. Each prediction was written before its run (Appendix~\ref{app:prereg}). WT/WB: write-through/write-back host tier; LFU/SLRU: least-frequently-used/segmented-LRU eviction; TTFT: time to first token.}
\label{fig:workflow}
\end{figure*}

